% BEGIN REV09_MULTIPLICIDAD_MODAL
\subsubsection{Multiplicity and recovery of modal itineraries}
\label{corpus9:modal:multiplicidad}

The incidence \(F_{\mathrm{av}}\), constructed by the tritic calendar
in the preceding development, acts here on its envelope of modal
walks. Write
\[
 F=F_{\mathrm{av}}=\begin{pmatrix}0&1\\1&1\end{pmatrix},
 \qquad 0=\Sigma,\quad 1=\Pi.
\]
The self-scaling coordinate \(\varphi\) and the positive ray
\(r=(1,\varphi)\) are the outputs already obtained from this
incidence. They determine the multiplicity of the scalar reading
and support a reversible representation of each itinerary.

Consider the set of finite directed walks
\(\gamma=(x_0,\ldots,x_n)\) of \(F\). Their endpoints and lengths
are distinguished, repetitions are allowed, and trivial walks are
included. This domain is the modal envelope; the complete TPK
state additionally retains the records of its enriched history.
The multiplicative character takes the form
\begin{equation}
 \chi(\gamma)=\varphi^n\frac{r_{x_0}}{r_{x_n}}
             =\varphi^{k(\gamma)},\qquad
 r=(1,\varphi),\quad k(\gamma)=n+x_0-x_n.
 \label{corpus9:modal:caracter}
\end{equation}
The transitions \(\Sigma\to\Pi\), \(\Pi\to\Pi\), and
\(\Pi\to\Sigma\) increment \(k\) by \(0\), \(1\), and \(2\),
respectively. The first step modifies the path while preserving
its character. In particular, both \((1,1)\) and \((0,1,1)\)
produce \(\varphi\) under this reading.

\begin{proposition}[Multiplicity of the self-scaling reading]
\label{corpus9:modal:conteo}
Let \(\mathcal F_k=\{\gamma:k(\gamma)=k\}\) and
\(N_k=|\mathcal F_k|\). Then
\begin{equation}
 N_0=3,\qquad
 N_k=2\operatorname{tr}(F^k)
     =2\bigl(\varphi^k+(-\varphi^{-1})^k\bigr)
     \quad(k\geq1),
 \label{corpus9:modal:multiplicidad-exacta}
\end{equation}
and
\[
 \lim_{k\to\infty}N_k^{1/k}=\varphi.
\]
\end{proposition}

\begin{proof}
Given endpoints \(i,j\in\{0,1\}\), the length is determined by
\(n=k-i+j\). The number of walks with those endpoints is
\((F^n)_{ij}\). For \(k\geq1\), summing the four sectors gives
\[
 N_k=\operatorname{tr}(F^k)
      +(F^{k+1})_{01}+(F^{k-1})_{10}.
\]
The identity \(F^2=F+I\) implies by induction that, for \(n\geq1\),
\[
 F^n=\begin{pmatrix}f_{n-1}&f_n\\f_n&f_{n+1}\end{pmatrix},
 \qquad f_0=0,\quad f_1=1,\quad f_{n+1}=f_n+f_{n-1}.
\]
For \(k\geq2\), the two off-diagonal terms are
\(f_{k+1}+f_{k-1}=\operatorname{tr}(F^k)\); for \(k=1\),
their sum is \(1+0=1=\operatorname{tr}(F)\), using \(F^0=I\).
For \(k=0\), the only walks are \((0)\), \((1)\), and \((0,1)\).
The eigenvalues of \(F\) are \(\varphi\) and \(-\varphi^{-1}\),
so its trace gives the closed formula. Finally,
\[
 N_k=2\varphi^k\bigl(1+(-\varphi^{-2})^k\bigr),
\]
and the factor in parentheses tends to one; taking \(k\)-th roots
yields the limit.
\end{proof}

Self-scaling thus coincides with the exponential growth rate
of the finite modal genealogies that the scalar character
makes indistinguishable.

\paragraph{Refinement by the composition of returns.}
Let \(b(\gamma)\) count the edges \(\Pi\to\Sigma\), and let
\(c(\gamma)\) count the edges \(\Pi\to\Pi\). The increments in
\eqref{corpus9:modal:caracter} give \(k=2b+c\).
For formal indeterminates \(z,t\), the graded incidence
\[
 M(z,t)=\begin{pmatrix}0&1\\tz^2&z\end{pmatrix}
\]
records exponent and returns jointly. Its only degree-zero edge
is \(\Sigma\to\Pi\), and there is no degree-zero cycle.
Each coefficient therefore counts a finite set of walks.
If \(N_{k,b}\) denotes their multiplicity with values \(k,b\),
the sum over all lengths satisfies
\begin{equation}
 \sum_{k,b\geq0}N_{k,b}z^kt^b
 =\mathbf1^{\mathsf T}(I-M(z,t))^{-1}\mathbf1
 =\frac{3-z+tz^2}{1-z-tz^2},
 \qquad \mathbf1=(1,1)^{\mathsf T}.
 \label{corpus9:modal:serie-graduada}
\end{equation}
Indeed, entry \(ij\) of \(M^n\) counts walks of length \(n\)
with those endpoints and their weights. Coefficientwise formal
summation yields \((I-M)^{-1}\). The determinant of \(I-M\)
is \(1-z-tz^2\), and the sum of the four entries of its adjugate
is \(3-z+tz^2\).

For \(k\geq1\) and \(0\leq b\leq\lfloor k/2\rfloor\), this gives
\begin{equation}
 N_{k,b}=\frac{2k}{k-b}\binom{k-b}{b}.
 \label{corpus9:modal:conteo-retornos}
\end{equation}
To prove this, the expansion
\[
 (1-z-tz^2)^{-1}=\sum_{m\geq0}(z+tz^2)^m
\]
gives coefficient \(\binom{k-b}{b}\) at \(z^kt^b\).
Multiplication by the numerator of
\eqref{corpus9:modal:serie-graduada} produces
\[
 3\binom{k-b}{b}-\binom{k-b-1}{b}
   +\binom{k-b-1}{b-1}
 =2\binom{k-b}{b}+2\binom{k-b-1}{b-1}.
\]
Coefficients outside their range are taken to be zero.
For \(b\geq1\), the identity
\(\binom{k-b-1}{b-1}=\frac{b}{k-b}\binom{k-b}{b}\)
yields \eqref{corpus9:modal:conteo-retornos}; for \(b=0\),
both sides equal two. The constant term is treated separately:
\(N_{0,0}=3\). For \(k=4\), there are fourteen walks, distributed
as \(2\), \(8\), and \(4\) according to whether they contain zero,
one, or two returns. The same character therefore admits distinct
internal compositions.

\begin{proposition}[Reversible recovery and transport of codes]
\label{corpus9:modal:recuperador}
Each walk in the modal envelope admits a code
\[
 \mathcal C(\gamma)=(k,a),\qquad 0\leq a<N_k,
\]
with a decoding map \(\mathcal D\) such that
\[
 \mathcal D\mathcal C=\mathrm{id},\qquad
 \mathcal C\mathcal D=\mathrm{id}.
\]
If \(A_v(\gamma)=\gamma v\) prolongs a walk by an admissible edge
and \(T\) removes the last state of a positive-length walk,
their representations
\[
 \mathcal E_v=\mathcal C A_v\mathcal D,\qquad
 \widehat T=\mathcal C T\mathcal D
\]
satisfy, on their respective domains,
\begin{equation}
 \widehat T\mathcal E_v=\mathrm{id},\qquad
 k(A_v\gamma)-k(\gamma)=1+x_n-v.
 \label{corpus9:modal:transporte-codigos}
\end{equation}
\end{proposition}

\begin{proof}
For each \(k\), the nonempty sectors are ordered by \((n,i,j)\)
and, within each sector, walks are ordered lexicographically by
their states. The number of length-\(m\) continuations from
\(v\) to \(j\) is \((F^m)_{vj}\). Traversing a walk, one sums
the sizes of preceding sectors and, at each position, the sizes
of the continuations lexicographically preceding the chosen one.
The result is its zero-based index \(a\).

To invert this process, the size of each preceding sector is
subtracted from the index until the unique containing sector is
identified. At each position, the same procedure is applied
to continuation sizes. These sets form a disjoint and exhaustive
partition of the starting sector. At length zero, the only walk
of a nonempty sector is its vertex. Assuming recovery of
continuations of length \(m-1\), the index intervals of each first
successor uniquely identify that successor and the remaining
index of its continuation. Induction on \(m\) proves that the
procedures are inverse, both on walks and on the indices
\(0,\ldots,N_k-1\).

Lexicographic order serializes the existing path; the operations
generating the path precede this serialization.
Since \(T A_v=\mathrm{id}\), the two inverse identities give
\[
 \widehat T\mathcal E_v
 =\mathcal C T\mathcal D\mathcal C A_v\mathcal D
 =\mathcal C T A_v\mathcal D
 =\mathrm{id}.
\]
Finally, the exponent definition gives
\[
 k(A_v\gamma)=(n+1)+x_0-v,\qquad
 k(\gamma)=n+x_0-x_n,
\]
whose difference is the second identity.
\end{proof}

\paragraph{Distinguishing capacity with a known exponent.}
With \(k\) known, an injective fixed-length binary representation
of length \(\ell\) requires \(2^\ell\geq N_k\), by the
pigeonhole principle. Consequently,
\(\ell\geq\lceil\log_2N_k\rceil\).
The preceding index \(a\) attains this length through its binary
representation, padded with leading zeros where necessary.
This minimum concerns transmitting the index with \(k\) known;
transmitting \(k\) and delimiting variable-length codes have
their own costs. Moreover,
\[
 \lim_{k\to\infty}\frac{\log_2N_k}{k}=\log_2\varphi.
\]
The equality measures the combinatorial capacity to distinguish
paths in the modal envelope.

\paragraph{Fiber information under an explicit distribution.}
For fixed \(k\), let \(R_k\) be a walk chosen uniformly from
\(\mathcal F_k\). We use natural logarithms and
\(H(R_k)=-\sum_{\gamma\in\mathcal F_k}p_\gamma\log p_\gamma\),
with \(p_\gamma=1/N_k\). All these walks have character
\(\varphi^k\), so
\[
 H(R_k\mid\chi(R_k))=\log N_k,\qquad
 H(R_k\mid\mathcal C(R_k))=0.
\]
The first equality follows because the reading is constant on
the fiber; the second follows from the unique decoding already
proved. In particular, \(\log N_k/k\to\log\varphi\).

For \(k=4\), the character leaves an uncertainty of \(\log14\)
nats. The three classes \(B=b(R_4)\) have probabilities
\(2/14\), \(8/14\), and \(4/14\). Conditional on each class,
the distribution remains uniform on its \(2\), \(8\), or \(4\)
elements, respectively. By the definition of conditional entropy,
\[
 H(R_4\mid B)
 =\frac2{14}\log2+\frac8{14}\log8+\frac4{14}\log4.
\]
The index within the class, together with \(B\), uniquely recovers
the walk. These values pertain to the declared uniform distribution
on the modal fiber. The measure on the complete space of histories,
the temporal realization of \(k\), and any thermal reading retain
their respective domains and transducers.
% END REV09_MULTIPLICIDAD_MODAL
